%Daten zur Institution

%\input{Dozentdaten}

%\renewcommand{\fachbereich}{Fachbereich}

%\renewcommand{\dozent}{Prof. Dr. . }

%Klausurdaten

\renewcommand{\klausurgebiet}{ }

\renewcommand{\klausurtyp}{ }

\renewcommand{\klausurdatum}{ . 20}

\klausurvorspann {\fachbereich} {\klausurdatum} {\dozent} {\klausurgebiet} {\klausurtyp}

%Daten für folgende Punktetabelle

\renewcommand{\aeins}{  3  }

\renewcommand{\azwei}{  3   }

\renewcommand{\adrei}{  0  }

\renewcommand{\avier}{  0  }

\renewcommand{\afuenf}{  4  }

\renewcommand{\asechs}{  5  }

\renewcommand{\asieben}{  4  }

\renewcommand{\aacht}{  0  }

\renewcommand{\aneun}{  4  }

\renewcommand{\azehn}{  6  }

\renewcommand{\aelf}{  2  }

\renewcommand{\azwoelf}{  6  }

\renewcommand{\adreizehn}{  4   }

\renewcommand{\avierzehn}{  10  }

\renewcommand{\afuenfzehn}{  10  }

\renewcommand{\asechzehn}{  0  }

\renewcommand{\asiebzehn}{  61  }

\renewcommand{\aachtzehn}{    }

\renewcommand{\aneunzehn}{    }

\renewcommand{\azwanzig}{    }

\renewcommand{\aeinundzwanzig}{    }

\renewcommand{\azweiundzwanzig}{    }

\renewcommand{\adreiundzwanzig}{    }

\renewcommand{\avierundzwanzig}{    }

\renewcommand{\afuenfundzwanzig}{    }

\renewcommand{\asechsundzwanzig}{    }

\punktetabellesechzehn

\klausurnote

\newpage

\setcounter{section}{K}

\inputaufgabegibtloesung
{3}
{

Definiere die folgenden
\zusatzklammer {kursiv gedruckten} {} {} Begriffe.
\aufzaehlungsechs{Die
\stichwort {Weingartenabbildung} {}
in $T_PY$ zu einem Punkt

\mavergleichskette
{\vergleichskette
{ P
}
{ \in }{ Y
}
{  }{
}
{    }{
}
{   }{
}
}
{}{}{}
auf einer orientierten differenzierbaren Hyperfläche

\mavergleichskette
{\vergleichskette
{ Y
}
{ \subseteq }{ \R^n
}
{  }{
}
{    }{
}
{   }{
}
}
{}{}{.}

}{Eine
\stichwort {reelles Vektorbündel} {}
$V$ vom Rang $r$ über einem
\definitionsverweis {topologischen Raum}{}{}
$X$.

}{Das \stichwort {Wegintegral} {} zu einer $1$-Differentialform
\mathl{\omega \in
{ \mathcal E }^{  1 } (  M   )}{} auf einer differenzierbaren Mannigfaltigkeit $M$ bezüglich einer stetig differenzierbaren Kurve
\maabb {\gamma} {[a,b]} {M
} {.}

}{Eine
\stichwort {positive Volumenform} {}
$\omega$ auf einer differenzierbaren Mannigfaltigkeit $M$.

}{Eine \stichwort {geschlossene} {} Differentialform $\omega$ auf einer differenzierbaren Mannigfaltigkeit $M$.

}{Die
\stichwort {tangentiale Beschleunigung} {}
einer zweifach stetig differenzierbaren Kurve
\maabb {\gamma} {I} {M
} {}
auf einer
\definitionsverweis {riemannschen Mannigfaltigkeit}{}{}
$M$.
}

}
{} {}

\inputaufgabegibtloesung
{3}
{

Formuliere die folgenden Sätze.
\aufzaehlungdrei{Der
\stichwort {Satz über die Normalkrümmung} {.}}{Die Formel für die Berechnung des kanonischen Volumens auf einer riemannschen Mannigfaltigkeit in einer Karte.}{Der \stichwort {Satz von Stokes} {} für Mannigfaltigkeiten mit Rand.}

}
{} {}

\inputaufgabe
{0}
{

}
{} {}

\inputaufgabe
{0}
{

}
{} {}

\inputaufgabegibtloesung
{4}
{

Beweise den Satz über die Beziehung zwischen Krümmung und Einheitsnormalenvektor einer bogenparametrisierten ebenen Kurve.

}
{} {}

\inputaufgabegibtloesung
{5 (1+2+2)}
{

Wir betrachten auf der zweidimensionalen
\definitionsverweis {Einheitssphäre}{}{}
$Y$ das
\definitionsverweis {tangentiale Vektorfeld}{}{}

\mavergleichskettedisp
{\vergleichskette
{ F \begin{pmatrix}  x  \\ y \\ z   \end{pmatrix}
}
{ =} { \begin{pmatrix}  y  \\ -x\\  0   \end{pmatrix}
}
{ } {
}
{ } {
}
{ } {
}
}
{}{}{.}
\aufzaehlungdrei{Bestimme die
\definitionsverweis {Jacobi-Matrix}{}{}
zu $F$ auf dem $\R^3$.
}{Bestimme für den Punkt

\mavergleichskette
{\vergleichskette
{ P
}
{ = }{ \begin{pmatrix}  0  \\ 1 \\  0   \end{pmatrix}
}
{ \in }{ Y
}
{    }{
}
{   }{
}
}
{}{}{}
eine Matrix für die Abbildung
\maabbeledisp {} {T_PY} { T_PY
} { v } { \nabla_v F
} {,}
bezüglich einer geeigneten
\definitionsverweis {Basis}{}{.}
}{Bestimme für den Punkt

\mavergleichskette
{\vergleichskette
{ Q
}
{ = }{ \begin{pmatrix}  0  \\ 0\\  1   \end{pmatrix}
}
{ \in }{ Y
}
{    }{
}
{   }{
}
}
{}{}{}
eine Matrix für die Abbildung
\maabbeledisp {} {T_QY} { T_QY
} { v } { \nabla_v F
} {,}
bezüglich einer geeigneten
\definitionsverweis {Basis}{}{.}
}

}
{} {}

\inputaufgabegibtloesung
{4 (2+2)}
{

Wir betrachten die Ellipse

\mavergleichskettedisp
{\vergleichskette
{ E
}
{ =} { \left\{ (u,v)  \mid   2u^2 +5v^2 = 4         \right\}
}
{ } {
}
{ } {
}
{ } {
}
}
{}{}{.}
\aufzaehlungdrei{Definiere eine surjektive stetig differenzierbare Abbildung
\maabb {} {\R} {E
} {.}
}{Beschreibe einen
\definitionsverweis {Diffeomorphismus}{}{}
zwischen der Sphäre $S^1$ und  $E$.
}{
}

}
{} {}

\inputaufgabe
{0}
{

}
{} {}

\inputaufgabegibtloesung
{4}
{

Man gebe ein
\definitionsverweis {stetiges}{}{} \definitionsverweis {Vektorfeld}{}{}
auf $S^2$ an, das nur eine Nullstelle besitzt.

}
{} {}

\inputaufgabegibtloesung
{6}
{

Es sei ein
\definitionsverweis {Verklebungsdatum}{}{}

\mathbed {U_i} {}
 {i \in I} {}
 {} {} {} {,}
für
\definitionsverweis {topologische Räume}{}{}
gegeben. Zeige, dass es einen eindeutig bestimmten topologischen Raum $X$, eine offene Überdeckung

\mavergleichskette
{\vergleichskette
{ X
}
{ = }{ \bigcup_{i \in I} V_i
}
{  }{
}
{    }{
}
{   }{
}
}
{}{}{}
und
\definitionsverweis {Homöomorphismen}{}{}
\maabb {\psi_i} {U_i } { V_i
} {}
derart gibt, dass

\mavergleichskettedisp
{\vergleichskette
{ \psi_i  \left(  U_{ij}  \right)
}
{ =} { V_i \cap V_j
}
{ } {
}
{ } {
}
{ } {
}
}
{}{}{}
ist und

\mavergleichskettedisp
{\vergleichskette
{ \psi_i \vert_{U_{ij} }
}
{ =} { \psi_j \vert_{U_{ji} }  \circ   \varphi_{ji}
}
{ } {
}
{ } {
}
{ } {
}
}
{}{}{}
gilt.

}
{} {}

\inputaufgabegibtloesung
{2}
{

Es sei $M$ eine
\definitionsverweis {differenzierbare Mannigfaltigkeit}{}{,}
\maabb {\gamma} { I } { M
} {}
eine
\definitionsverweis {differenzierbare Kurve}{}{}
in $M$ und $\omega$ eine differenzierbare
$1$-\definitionsverweis {Form}{}{}
auf $M$. Zeige, dass

\mavergleichskettedisp
{\vergleichskette
{ \gamma^* (d \omega)
}
{ =} { 0
}
{ } {
}
{ } {
}
{ } {
}
}
{}{}{}
ist.

}
{} {}

\inputaufgabegibtloesung
{6 (3+3)}
{

Es sei $M$ eine
\definitionsverweis {kompakte}{}{}
\definitionsverweis {differenzierbare Mannigfaltigkeit}{}{}
mit zumindest zwei Punkten.
\aufzaehlungzwei {Zeige, dass eine
\definitionsverweis {exakte}{}{}
stetige
$1$-\definitionsverweis {Differentialform}{}{}
auf $M$ zumindest zwei Nullstellen besitzt.
} {Zeige, dass dies für eine
\definitionsverweis {geschlossene}{}{}
Differentialform nicht gelten muss.
}

}
{} {}

\inputaufgabegibtloesung
{4}
{

Man gebe ein Beispiel für eine
\definitionsverweis {zusammenhängende}{}{}
\definitionsverweis {Mannigfaltigkeit mit Rand}{}{,}
deren Rand aus
\zusatzklammer {einer disjunkten Vereinigung von} {} {}
unendlich vielen  $\R^1$ besteht.

}
{} {}

\inputaufgabegibtloesung
{10}
{

Beweise den  \stichwort {Satz über die Partition der Eins} {.}

}
{} {}

\inputaufgabegibtloesung
{10 (2+6+2)}
{

Wir betrachten die
$2$-\definitionsverweis {Differentialform}{}{}

\mavergleichskettedisp
{\vergleichskette
{ \omega
}
{ =} { x dy \wedge dz -y dx \wedge dz + z dx \wedge dy
}
{ } {
}
{ } {
}
{ } {
}
}
{}{}{}
auf der Einheitskugel
\mathl{B  \left(  0 , 1  \right)}{.}
\aufzaehlungdreiabc{Zeige, dass $d \omega$ das Dreifache der Standardvolumenform auf der Kugel ist.
}{Zeige, dass
\mathl{\omega\vert_{S^2}}{} die Standardflächenform auf der Einheitssphäre ist.
}{Berechne die Kugeloberfläche aus dem
[[Allgemeines Kugelvolumen/Mit Cavalieri-Prinzip/Beispiel|Kugelvolumen]]
mit dem
[[Mannigfaltigkeit mit Rand/Satz von Stokes/Fakt|Satz von Stokes]].
}

}
{} {}

\inputaufgabe
{0}
{

}
{} {}

 [[Kategorie:Latexseite]]
